% Part: many-valued-logic
% Chapter: syntax-and-semantics
% Section: connectives

\documentclass[../../../include/open-logic-section]{subfiles}

\begin{document}

\olfileid{mvl}{syn}{con}

\olsection{భాషలు, సంయోజకాలు}

సాంప్రదాయిక ప్రతిజ్ఞావాక్య తర్కం, ఇంకా అనేక ఇతర తర్కాలు,
\emph{ప్రతిజ్ఞావాక్య స్థిరాంకాలు}, \emph{సంయోజకాలు} గల ఒక
నిర్దిష్ట సముదాయాన్ని వాడతాయి. ఉదాహరణకు, కింది వాటిని
ప్రాథమిక సంకేతాలుగా తీసుకుంటాం:
\begin{enumerate}
\tagitem{prvFalse}{\tetoken{అసత్యం}{!!{falsity}} కోసం ప్రతిజ్ఞావాక్య స్థిరాంకం~$\lfalse$.}{}
\tagitem{prvTrue}{\tetoken{సత్యం}{!!{truth}} కోసం ప్రతిజ్ఞావాక్య స్థిరాంకం~$\ltrue$.}{}
\item తార్కిక సంయోజకాలు:
  \startycommalist
  \iftag{prvNot}{\ycomma $\lnot$ (నిషేధం)}{}%
  \iftag{prvAnd}{\ycomma $\land$ (సంయోగం)}{}%
  \iftag{prvOr}{\ycomma $\lor$ (వికల్పం)}{}%
  \iftag{prvIf}{\ycomma $\lif$ (\tetoken{సోపాధికం}{!!{conditional}})}{}%
  \iftag{prvIff}{\ycomma $\liff$ (\tetoken{ద్విసోపాధికం}{!!{biconditional}})}{}%
\end{enumerate}
\iftag{defNot,defOr,defAnd,defIf,defIff,defTrue,defFalse,defEx,defAll}{%
పైన ఉన్న ప్రాథమిక సంయోజకాలతో పాటు, సంక్షిప్తీకరణలుగా
నిర్వచించిన కింది సంకేతాలనూ వాడతాం:
\startycommalist
  \iftag{defNot}{\ycomma $\lnot$ (నిషేధం)}{}%
  \iftag{defAnd}{\ycomma $\land$ (సంయోగం)}{}%
  \iftag{defOr}{\ycomma $\lor$ (వికల్పం)}{}%
  \iftag{defIf}{\ycomma $\lif$ (\tetoken{సోపాధికం}{!!{conditional}})}{}%
  \iftag{defIff}{\ycomma $\liff$ (\tetoken{ద్విసోపాధికం}{!!{biconditional}})}{}%
  \iftag{defFalse}{\ycomma $\lfalse$ (\tetoken{అసత్యం}{!!{falsity}})}{}%
  \iftag{defTrue}{\ycomma $\ltrue$ (\tetoken{సత్యం}{!!{truth}})}{}.
}{}
\sourcecorrection{OLTEMVLSYN-001}{నిర్వచిత సంకేతాల జాబితా చివర ఆంగ్ల మూలంలో
లోపలి \texttt{defTrue} పరీక్షకు ఖాళీ అసత్య శాఖ లేకుండానే బయటి పరీక్ష
ముగుస్తుంది. లోపలి పరీక్షను రెండు శాఖలతో ముగించి, తరువాత బయటి పరీక్షను
ముగిసేలా కుండలీకరణల క్రమాన్ని సరిచేశాం; సంకేతాల జాబితా మారలేదు.}

బహుమూల్య తర్కాల్లోనూ ఇవే సంయోజకాలను వాడతారు. అయితే ఒకే తర్కంలో,
ఉదాహరణకు సంయోగానికి, వేర్వేరు రూపాలను చేర్చడం తరచుగా ఉపయోగకరం;
వాటిని వేరు చూపడానికి వేర్వేరు సంకేతాలు కావాలి. కొన్ని బహుమూల్య
తర్కాల్లో సాంప్రదాయిక తర్కంలో సమానార్థకం లేని సంయోజకాలు కూడా ఉంటాయి.
కాబట్టి ఇక్కడ సాధారణ పద్ధతి కంటే కొంత విస్తృతంగా వ్యవహరిస్తాం.

\begin{defn}
  \emph{ప్రతిజ్ఞావాక్య భాష} అంటే సంయోజకాల సమితి~$\Lang L$.
  ప్రతి సంయోజకం~$\star$కు ఒక \emph{స్థానసంఖ్య} ఉంటుంది;
  స్థానసంఖ్య~$n$ గల సంయోజకాన్ని \emph{$n$-స్థానికం} అంటాం.
  స్థానసంఖ్య~$0$ గల సంయోజకాలను \emph{స్థిరాంకాలు} అని కూడా,
  స్థానసంఖ్య~$1$ గల వాటిని \emph{ఏకస్థానికాలు} అని,
  స్థానసంఖ్య~$2$ గల వాటిని \emph{ద్విస్థానికాలు} అని అంటాం.
\end{defn}

\begin{ex}
  ప్రతిజ్ఞావాక్య తర్కపు ప్రామాణిక భాష~$\Lang L_0$లో కింది
  సంయోజకాలు ఉంటాయి (వాటి స్థానసంఖ్యలు కుండలీకరణల్లో):
  $\lfalse$~($0$),
  $\lnot$~($1$),
  $\land$~($2$),
  $\lor$~($2$),
  $\lif$~($2$). మనం పరిశీలించే చాలా తర్కాలు ఈ భాషను వాడతాయి.
  కొన్ని తర్కాలు సంప్రదాయానుసారంగా కొన్ని సంయోజకాలకు వేరు
  సంకేతాలను వాడతాయి. ఉదాహరణకు, గుణిత తర్కంలో సంయోగానికి
  $\land$ బదులు తరచుగా~$\odot$ వాడతారు. కొన్నిసార్లు కొత్త
  సంచాలకాన్ని చేర్చడం అనుకూలంగా ఉంటుంది; ఉదాహరణకు,
  నిర్ణీతత్వ సంచాలకం~$\triangle$ ($1$-స్థానికం).
\end{ex}

\end{document}
